\subsection*{Textbook Formal Inventory}
This subsection records the exact formal material in \file{HAETAE_Assumptions.ec}.  It is generated from the proof-surface EasyCrypt file, so the line numbers and statements are meant to be read next to the checked source.

\noindent\textbf{Chapter role.} formal MLWE, Module-SIS, and bimodal self-target MSIS games.

\noindent\textbf{Inventory size.} 11 context declarations, 60 definitions/game objects, 11 lemmas or relational theorems, and 0 explicit Hoare/Phoare judgments were found.

\subsubsection*{Theory Context, Imports, and Quantified Modules}
\paragraph{Context 1 (line 1)}
\noindent\textbf{EasyCrypt name.}
\begin{lstlisting}[language=EasyCrypt,basicstyle=\ttfamily\scriptsize]
imported theories
\end{lstlisting}
\noindent\textbf{Checked EasyCrypt statement.}
\begin{lstlisting}[language=EasyCrypt,basicstyle=\ttfamily\scriptsize]
require import AllCore Distr List.
\end{lstlisting}
\noindent\textbf{Formal mathematical reading.}
Mathematical context. This line imports or specializes earlier theories, making their carrier sets, functions, modules, and theorems available to the current file.

\noindent\textbf{Role in the proof.}
Use in this chapter. It fixes the proof context, imported mathematical language, or quantified module parameters for the following statements.

\paragraph{Context 2 (line 2)}
\noindent\textbf{EasyCrypt name.}
\begin{lstlisting}[language=EasyCrypt,basicstyle=\ttfamily\scriptsize]
imported theories
\end{lstlisting}
\noindent\textbf{Checked EasyCrypt statement.}
\begin{lstlisting}[language=EasyCrypt,basicstyle=\ttfamily\scriptsize]
require import HAETAE_Params HAETAE_Algebra.
\end{lstlisting}
\noindent\textbf{Formal mathematical reading.}
Mathematical context. This line imports or specializes earlier theories, making their carrier sets, functions, modules, and theorems available to the current file.

\noindent\textbf{Role in the proof.}
Use in this chapter. It fixes the proof context, imported mathematical language, or quantified module parameters for the following statements.

\paragraph{Context 3 (line 4)}
\noindent\textbf{EasyCrypt name.}
\begin{lstlisting}[language=EasyCrypt,basicstyle=\ttfamily\scriptsize]
HAETAE_Assumptions
\end{lstlisting}
\noindent\textbf{Checked EasyCrypt statement.}
\begin{lstlisting}[language=EasyCrypt,basicstyle=\ttfamily\scriptsize]
theory HAETAE_Assumptions.
\end{lstlisting}
\noindent\textbf{Formal mathematical reading.}
Mathematical context. This begins a namespace for the formal development in this file.

\noindent\textbf{Role in the proof.}
Use in this chapter. It fixes the proof context, imported mathematical language, or quantified module parameters for the following statements.

\paragraph{Context 4 (line 6)}
\noindent\textbf{EasyCrypt name.}
\begin{lstlisting}[language=EasyCrypt,basicstyle=\ttfamily\scriptsize]
opened namespace
\end{lstlisting}
\noindent\textbf{Checked EasyCrypt statement.}
\begin{lstlisting}[language=EasyCrypt,basicstyle=\ttfamily\scriptsize]
import HAETAE_Params.
\end{lstlisting}
\noindent\textbf{Formal mathematical reading.}
Mathematical context. This line imports or specializes earlier theories, making their carrier sets, functions, modules, and theorems available to the current file.

\noindent\textbf{Role in the proof.}
Use in this chapter. It fixes the proof context, imported mathematical language, or quantified module parameters for the following statements.

\paragraph{Context 5 (line 7)}
\noindent\textbf{EasyCrypt name.}
\begin{lstlisting}[language=EasyCrypt,basicstyle=\ttfamily\scriptsize]
opened namespace
\end{lstlisting}
\noindent\textbf{Checked EasyCrypt statement.}
\begin{lstlisting}[language=EasyCrypt,basicstyle=\ttfamily\scriptsize]
import HAETAE_Algebra.
\end{lstlisting}
\noindent\textbf{Formal mathematical reading.}
Mathematical context. This line imports or specializes earlier theories, making their carrier sets, functions, modules, and theorems available to the current file.

\noindent\textbf{Role in the proof.}
Use in this chapter. It fixes the proof context, imported mathematical language, or quantified module parameters for the following statements.

\paragraph{Context 6 (line 205)}
\noindent\textbf{EasyCrypt name.}
\begin{lstlisting}[language=EasyCrypt,basicstyle=\ttfamily\scriptsize]
BimodalSolverLift
\end{lstlisting}
\noindent\textbf{Checked EasyCrypt statement.}
\begin{lstlisting}[language=EasyCrypt,basicstyle=\ttfamily\scriptsize]
section BimodalSolverLift.
\end{lstlisting}
\noindent\textbf{Formal mathematical reading.}
Mathematical context. This opens a scoped block of hypotheses, module parameters, and lemmas.

\noindent\textbf{Role in the proof.}
Use in this chapter. It fixes the proof context, imported mathematical language, or quantified module parameters for the following statements.

\paragraph{Context 7 (line 207)}
\noindent\textbf{EasyCrypt name.}
\begin{lstlisting}[language=EasyCrypt,basicstyle=\ttfamily\scriptsize]
B
\end{lstlisting}
\noindent\textbf{Checked EasyCrypt statement.}
\begin{lstlisting}[language=EasyCrypt,basicstyle=\ttfamily\scriptsize]
declare module B <: BimodalSelfTargetMSIS_Solver.
\end{lstlisting}
\noindent\textbf{Formal mathematical reading.}
Mathematical context. This is universal quantification over an adversary, oracle, scheme, sampler, or reduction module satisfying the stated interface and memory restrictions.

\noindent\textbf{Role in the proof.}
Use in this chapter. It fixes the proof context, imported mathematical language, or quantified module parameters for the following statements.

\paragraph{Context 8 (line 288)}
\noindent\textbf{EasyCrypt name.}
\begin{lstlisting}[language=EasyCrypt,basicstyle=\ttfamily\scriptsize]
BimodalModeSolverLiftedDistribution
\end{lstlisting}
\noindent\textbf{Checked EasyCrypt statement.}
\begin{lstlisting}[language=EasyCrypt,basicstyle=\ttfamily\scriptsize]
section BimodalModeSolverLiftedDistribution.
\end{lstlisting}
\noindent\textbf{Formal mathematical reading.}
Mathematical context. This opens a scoped block of hypotheses, module parameters, and lemmas.

\noindent\textbf{Role in the proof.}
Use in this chapter. It fixes the proof context, imported mathematical language, or quantified module parameters for the following statements.

\paragraph{Context 9 (line 290)}
\noindent\textbf{EasyCrypt name.}
\begin{lstlisting}[language=EasyCrypt,basicstyle=\ttfamily\scriptsize]
B
\end{lstlisting}
\noindent\textbf{Checked EasyCrypt statement.}
\begin{lstlisting}[language=EasyCrypt,basicstyle=\ttfamily\scriptsize]
declare module B <: BimodalSelfTargetMSIS_ModeSolver.
\end{lstlisting}
\noindent\textbf{Formal mathematical reading.}
Mathematical context. This is universal quantification over an adversary, oracle, scheme, sampler, or reduction module satisfying the stated interface and memory restrictions.

\noindent\textbf{Role in the proof.}
Use in this chapter. It fixes the proof context, imported mathematical language, or quantified module parameters for the following statements.

\paragraph{Context 10 (line 358)}
\noindent\textbf{EasyCrypt name.}
\begin{lstlisting}[language=EasyCrypt,basicstyle=\ttfamily\scriptsize]
BimodalModuleSISReduction
\end{lstlisting}
\noindent\textbf{Checked EasyCrypt statement.}
\begin{lstlisting}[language=EasyCrypt,basicstyle=\ttfamily\scriptsize]
section BimodalModuleSISReduction.
\end{lstlisting}
\noindent\textbf{Formal mathematical reading.}
Mathematical context. This opens a scoped block of hypotheses, module parameters, and lemmas.

\noindent\textbf{Role in the proof.}
Use in this chapter. It fixes the proof context, imported mathematical language, or quantified module parameters for the following statements.

\paragraph{Context 11 (line 360)}
\noindent\textbf{EasyCrypt name.}
\begin{lstlisting}[language=EasyCrypt,basicstyle=\ttfamily\scriptsize]
B
\end{lstlisting}
\noindent\textbf{Checked EasyCrypt statement.}
\begin{lstlisting}[language=EasyCrypt,basicstyle=\ttfamily\scriptsize]
declare module B <: BimodalSelfTargetMSIS_Reduction.
\end{lstlisting}
\noindent\textbf{Formal mathematical reading.}
Mathematical context. This is universal quantification over an adversary, oracle, scheme, sampler, or reduction module satisfying the stated interface and memory restrictions.

\noindent\textbf{Role in the proof.}
Use in this chapter. It fixes the proof context, imported mathematical language, or quantified module parameters for the following statements.

\subsubsection*{Definitions, Carrier Sets, Operations, Games, and Procedures}
\begin{formaldefinition}[EasyCrypt line 9]
\noindent\textbf{EasyCrypt name.}
\begin{lstlisting}[language=EasyCrypt,basicstyle=\ttfamily\scriptsize]
mlwe_instance
\end{lstlisting}
\noindent\textbf{Checked EasyCrypt statement.}
\begin{lstlisting}[language=EasyCrypt,basicstyle=\ttfamily\scriptsize]
type mlwe_instance.
\end{lstlisting}
\noindent\textbf{Formal mathematical reading.}
Mathematical definition. \(\mathsf{mlwe\_instance}\) is an abstract carrier set.  The proof may quantify over this domain, but it does not expose a representation.

\noindent\textbf{Role in the proof.}
Use in this chapter. It is part of the exact formal vocabulary used by subsequent lemmas and games.

\end{formaldefinition}

\begin{formaldefinition}[EasyCrypt line 10]
\noindent\textbf{EasyCrypt name.}
\begin{lstlisting}[language=EasyCrypt,basicstyle=\ttfamily\scriptsize]
module_sis_instance
\end{lstlisting}
\noindent\textbf{Checked EasyCrypt statement.}
\begin{lstlisting}[language=EasyCrypt,basicstyle=\ttfamily\scriptsize]
type module_sis_instance = matrix * polyveck.
\end{lstlisting}
\noindent\textbf{Formal mathematical reading.}
Mathematical definition. This is a definitional type abbreviation.  The exact displayed EasyCrypt statement gives the right-hand side; later proof steps may unfold it without adding an assumption.

\noindent\textbf{Role in the proof.}
Use in this chapter. It is part of the exact formal vocabulary used by subsequent lemmas and games.

\end{formaldefinition}

\begin{formaldefinition}[EasyCrypt line 11]
\noindent\textbf{EasyCrypt name.}
\begin{lstlisting}[language=EasyCrypt,basicstyle=\ttfamily\scriptsize]
module_sis_solution
\end{lstlisting}
\noindent\textbf{Checked EasyCrypt statement.}
\begin{lstlisting}[language=EasyCrypt,basicstyle=\ttfamily\scriptsize]
type module_sis_solution = polyvecl.
\end{lstlisting}
\noindent\textbf{Formal mathematical reading.}
Mathematical definition. This is a definitional type abbreviation.  The exact displayed EasyCrypt statement gives the right-hand side; later proof steps may unfold it without adding an assumption.

\noindent\textbf{Role in the proof.}
Use in this chapter. It is part of the exact formal vocabulary used by subsequent lemmas and games.

\end{formaldefinition}

\begin{formaldefinition}[EasyCrypt line 12]
\noindent\textbf{EasyCrypt name.}
\begin{lstlisting}[language=EasyCrypt,basicstyle=\ttfamily\scriptsize]
bimodal_selftarget_msis_instance
\end{lstlisting}
\noindent\textbf{Checked EasyCrypt statement.}
\begin{lstlisting}[language=EasyCrypt,basicstyle=\ttfamily\scriptsize]
type bimodal_selftarget_msis_instance = module_sis_instance.
\end{lstlisting}
\noindent\textbf{Formal mathematical reading.}
Mathematical definition. This is a definitional type abbreviation.  The exact displayed EasyCrypt statement gives the right-hand side; later proof steps may unfold it without adding an assumption.

\noindent\textbf{Role in the proof.}
Use in this chapter. It is part of the exact formal vocabulary used by subsequent lemmas and games.

\end{formaldefinition}

\begin{formaldefinition}[EasyCrypt line 13]
\noindent\textbf{EasyCrypt name.}
\begin{lstlisting}[language=EasyCrypt,basicstyle=\ttfamily\scriptsize]
bimodal_selftarget_msis_solution
\end{lstlisting}
\noindent\textbf{Checked EasyCrypt statement.}
\begin{lstlisting}[language=EasyCrypt,basicstyle=\ttfamily\scriptsize]
type bimodal_selftarget_msis_solution = module_sis_solution.
\end{lstlisting}
\noindent\textbf{Formal mathematical reading.}
Mathematical definition. This is a definitional type abbreviation.  The exact displayed EasyCrypt statement gives the right-hand side; later proof steps may unfold it without adding an assumption.

\noindent\textbf{Role in the proof.}
Use in this chapter. It is part of the exact formal vocabulary used by subsequent lemmas and games.

\end{formaldefinition}

\begin{formaldefinition}[EasyCrypt line 15]
\noindent\textbf{EasyCrypt name.}
\begin{lstlisting}[language=EasyCrypt,basicstyle=\ttfamily\scriptsize]
dmlwe_real
\end{lstlisting}
\noindent\textbf{Checked EasyCrypt statement.}
\begin{lstlisting}[language=EasyCrypt,basicstyle=\ttfamily\scriptsize]
op [lossless] dmlwe_real : mode -> mlwe_instance distr.
\end{lstlisting}
\noindent\textbf{Formal mathematical reading.}
Mathematical definition. This is a total EasyCrypt operation.  The exact displayed statement gives the domain, codomain, and defining equation when present.  The EasyCrypt [lossless] annotation states that this distribution has total mass 1.

\noindent\textbf{Role in the proof.}
Use in this chapter. It contributes a sampler or sampler-derived object to the distributional model.

\end{formaldefinition}

\begin{formaldefinition}[EasyCrypt line 16]
\noindent\textbf{EasyCrypt name.}
\begin{lstlisting}[language=EasyCrypt,basicstyle=\ttfamily\scriptsize]
dmlwe_random
\end{lstlisting}
\noindent\textbf{Checked EasyCrypt statement.}
\begin{lstlisting}[language=EasyCrypt,basicstyle=\ttfamily\scriptsize]
op [lossless] dmlwe_random : mode -> mlwe_instance distr.
\end{lstlisting}
\noindent\textbf{Formal mathematical reading.}
Mathematical definition. This is a total EasyCrypt operation.  The exact displayed statement gives the domain, codomain, and defining equation when present.  The EasyCrypt [lossless] annotation states that this distribution has total mass 1.

\noindent\textbf{Role in the proof.}
Use in this chapter. It contributes a sampler or sampler-derived object to the distributional model.

\end{formaldefinition}

\begin{formaldefinition}[EasyCrypt line 17]
\noindent\textbf{EasyCrypt name.}
\begin{lstlisting}[language=EasyCrypt,basicstyle=\ttfamily\scriptsize]
dmodule_sis_instance
\end{lstlisting}
\noindent\textbf{Checked EasyCrypt statement.}
\begin{lstlisting}[language=EasyCrypt,basicstyle=\ttfamily\scriptsize]
op [lossless] dmodule_sis_instance : mode -> module_sis_instance distr.
\end{lstlisting}
\noindent\textbf{Formal mathematical reading.}
Mathematical definition. This is a total EasyCrypt operation.  The exact displayed statement gives the domain, codomain, and defining equation when present.  The EasyCrypt [lossless] annotation states that this distribution has total mass 1.

\noindent\textbf{Role in the proof.}
Use in this chapter. It contributes a sampler or sampler-derived object to the distributional model.

\end{formaldefinition}

\begin{formaldefinition}[EasyCrypt line 18]
\noindent\textbf{EasyCrypt name.}
\begin{lstlisting}[language=EasyCrypt,basicstyle=\ttfamily\scriptsize]
dbimodal_selftarget_msis_instance
\end{lstlisting}
\noindent\textbf{Checked EasyCrypt statement.}
\begin{lstlisting}[language=EasyCrypt,basicstyle=\ttfamily\scriptsize]
op [lossless] dbimodal_selftarget_msis_instance :
  mode -> bimodal_selftarget_msis_instance distr.
\end{lstlisting}
\noindent\textbf{Formal mathematical reading.}
Mathematical definition. This is a total EasyCrypt operation.  The exact displayed statement gives the domain, codomain, and defining equation when present.  The EasyCrypt [lossless] annotation states that this distribution has total mass 1.

\noindent\textbf{Role in the proof.}
Use in this chapter. It contributes a sampler or sampler-derived object to the distributional model.

\end{formaldefinition}

\begin{formaldefinition}[EasyCrypt line 21]
\noindent\textbf{EasyCrypt name.}
\begin{lstlisting}[language=EasyCrypt,basicstyle=\ttfamily\scriptsize]
module_sis_eval
\end{lstlisting}
\noindent\textbf{Checked EasyCrypt statement.}
\begin{lstlisting}[language=EasyCrypt,basicstyle=\ttfamily\scriptsize]
op module_sis_eval
   (md : mode) (x : module_sis_instance)
   (sol : module_sis_solution) : polyveck =
  polyveck_sub md (matrix_vec_mul md x.`1 sol) x.`2.
\end{lstlisting}
\noindent\textbf{Formal mathematical reading.}
Mathematical definition. This is a total EasyCrypt operation.  The exact displayed statement gives the domain, codomain, and defining equation when present.

\noindent\textbf{Role in the proof.}
Use in this chapter. It is part of the exact formal vocabulary used by subsequent lemmas and games.

\end{formaldefinition}

\begin{formaldefinition}[EasyCrypt line 26]
\noindent\textbf{EasyCrypt name.}
\begin{lstlisting}[language=EasyCrypt,basicstyle=\ttfamily\scriptsize]
module_sis_zero_instance
\end{lstlisting}
\noindent\textbf{Checked EasyCrypt statement.}
\begin{lstlisting}[language=EasyCrypt,basicstyle=\ttfamily\scriptsize]
op module_sis_zero_instance (md : mode) : module_sis_instance =
  ([], polyveck_zero md).
\end{lstlisting}
\noindent\textbf{Formal mathematical reading.}
Mathematical definition. This is a total EasyCrypt operation.  The exact displayed statement gives the domain, codomain, and defining equation when present.

\noindent\textbf{Role in the proof.}
Use in this chapter. It is part of the exact formal vocabulary used by subsequent lemmas and games.

\end{formaldefinition}

\begin{formaldefinition}[EasyCrypt line 29]
\noindent\textbf{EasyCrypt name.}
\begin{lstlisting}[language=EasyCrypt,basicstyle=\ttfamily\scriptsize]
module_sis_solution_wf
\end{lstlisting}
\noindent\textbf{Checked EasyCrypt statement.}
\begin{lstlisting}[language=EasyCrypt,basicstyle=\ttfamily\scriptsize]
op module_sis_solution_wf (md : mode) (sol : module_sis_solution) : bool =
  polyvecl_wf md sol.
\end{lstlisting}
\noindent\textbf{Formal mathematical reading.}
Mathematical definition. This is a Boolean predicate.  The exact displayed EasyCrypt statement gives its arguments; mathematically it is an event, invariant, support condition, or well-formedness predicate over those arguments.

\noindent\textbf{Role in the proof.}
Use in this chapter. It names a well-formedness, support, or validity predicate needed by later preservation lemmas.

\end{formaldefinition}

\begin{formaldefinition}[EasyCrypt line 32]
\noindent\textbf{EasyCrypt name.}
\begin{lstlisting}[language=EasyCrypt,basicstyle=\ttfamily\scriptsize]
module_sis_valid
\end{lstlisting}
\noindent\textbf{Checked EasyCrypt statement.}
\begin{lstlisting}[language=EasyCrypt,basicstyle=\ttfamily\scriptsize]
op module_sis_valid
   (md : mode) (x : module_sis_instance)
   (sol : module_sis_solution) : bool =
  module_sis_solution_wf md sol /\
  matrix_vec_mul md x.`1 sol = x.`2.
\end{lstlisting}
\noindent\textbf{Formal mathematical reading.}
Mathematical definition. This is a Boolean predicate.  The exact displayed EasyCrypt statement gives its arguments; mathematically it is an event, invariant, support condition, or well-formedness predicate over those arguments.

\noindent\textbf{Role in the proof.}
Use in this chapter. It names a well-formedness, support, or validity predicate needed by later preservation lemmas.

\end{formaldefinition}

\begin{formaldefinition}[EasyCrypt line 38]
\noindent\textbf{EasyCrypt name.}
\begin{lstlisting}[language=EasyCrypt,basicstyle=\ttfamily\scriptsize]
bimodal_selftarget_msis_valid
\end{lstlisting}
\noindent\textbf{Checked EasyCrypt statement.}
\begin{lstlisting}[language=EasyCrypt,basicstyle=\ttfamily\scriptsize]
op bimodal_selftarget_msis_valid
   (md : mode) (x : bimodal_selftarget_msis_instance)
   (sol : bimodal_selftarget_msis_solution) : bool =
  module_sis_valid md x sol.
\end{lstlisting}
\noindent\textbf{Formal mathematical reading.}
Mathematical definition. This is a Boolean predicate.  The exact displayed EasyCrypt statement gives its arguments; mathematically it is an event, invariant, support condition, or well-formedness predicate over those arguments.

\noindent\textbf{Role in the proof.}
Use in this chapter. It names a well-formedness, support, or validity predicate needed by later preservation lemmas.

\end{formaldefinition}

\begin{formaldefinition}[EasyCrypt line 72]
\noindent\textbf{EasyCrypt name.}
\begin{lstlisting}[language=EasyCrypt,basicstyle=\ttfamily\scriptsize]
bimodal_to_module_sis_instance
\end{lstlisting}
\noindent\textbf{Checked EasyCrypt statement.}
\begin{lstlisting}[language=EasyCrypt,basicstyle=\ttfamily\scriptsize]
op bimodal_to_module_sis_instance :
  mode -> bimodal_selftarget_msis_instance -> module_sis_instance =
  fun (_ : mode) x => x.
\end{lstlisting}
\noindent\textbf{Formal mathematical reading.}
Mathematical definition. This is a total EasyCrypt operation.  The exact displayed statement gives the domain, codomain, and defining equation when present.

\noindent\textbf{Role in the proof.}
Use in this chapter. It is part of the exact formal vocabulary used by subsequent lemmas and games.

\end{formaldefinition}

\begin{formaldefinition}[EasyCrypt line 75]
\noindent\textbf{EasyCrypt name.}
\begin{lstlisting}[language=EasyCrypt,basicstyle=\ttfamily\scriptsize]
bimodal_to_module_sis_solution
\end{lstlisting}
\noindent\textbf{Checked EasyCrypt statement.}
\begin{lstlisting}[language=EasyCrypt,basicstyle=\ttfamily\scriptsize]
op bimodal_to_module_sis_solution :
  mode -> bimodal_selftarget_msis_instance ->
  bimodal_selftarget_msis_solution -> module_sis_solution =
  fun (_ : mode) (_ : bimodal_selftarget_msis_instance) sol => sol.
\end{lstlisting}
\noindent\textbf{Formal mathematical reading.}
Mathematical definition. This is a total EasyCrypt operation.  The exact displayed statement gives the domain, codomain, and defining equation when present.

\noindent\textbf{Role in the proof.}
Use in this chapter. It is part of the exact formal vocabulary used by subsequent lemmas and games.

\end{formaldefinition}

\begin{formaldefinition}[EasyCrypt line 91]
\noindent\textbf{EasyCrypt name.}
\begin{lstlisting}[language=EasyCrypt,basicstyle=\ttfamily\scriptsize]
MLWE_Distinguisher
\end{lstlisting}
\noindent\textbf{Checked EasyCrypt statement.}
\begin{lstlisting}[language=EasyCrypt,basicstyle=\ttfamily\scriptsize]
module type MLWE_Distinguisher = {
\end{lstlisting}
\noindent\textbf{Formal mathematical reading.}
Mathematical definition. This is an interface for an oracle, adversary, scheme, sampler, solver, or reduction.  A later theorem quantified over this interface is universally quantified over all modules implementing these procedures.

\noindent\textbf{Role in the proof.}
Use in this chapter. It supplies an executable component of the formal probabilistic model.

\end{formaldefinition}

\begin{formaldefinition}[EasyCrypt line 92]
\noindent\textbf{EasyCrypt name.}
\begin{lstlisting}[language=EasyCrypt,basicstyle=\ttfamily\scriptsize]
distinguish
\end{lstlisting}
\noindent\textbf{Checked EasyCrypt statement.}
\begin{lstlisting}[language=EasyCrypt,basicstyle=\ttfamily\scriptsize]
proc distinguish(x : mlwe_instance) : bool
}.
\end{lstlisting}
\noindent\textbf{Formal mathematical reading.}
Mathematical definition. This procedure is a command in the probabilistic program.  Its return value, state updates, and oracle calls are the operational content of the game.

\noindent\textbf{Role in the proof.}
Use in this chapter. It supplies an executable component of the formal probabilistic model.

\end{formaldefinition}

\begin{formaldefinition}[EasyCrypt line 95]
\noindent\textbf{EasyCrypt name.}
\begin{lstlisting}[language=EasyCrypt,basicstyle=\ttfamily\scriptsize]
MLWE_Real_Game
\end{lstlisting}
\noindent\textbf{Checked EasyCrypt statement.}
\begin{lstlisting}[language=EasyCrypt,basicstyle=\ttfamily\scriptsize]
module MLWE_Real_Game(B : MLWE_Distinguisher) = {
\end{lstlisting}
\noindent\textbf{Formal mathematical reading.}
Mathematical definition. This is a probabilistic program/game or a module adapter.  Its procedures induce the probability space used by the game-based proof.

\noindent\textbf{Role in the proof.}
Use in this chapter. It defines the game or game interface whose success event is bounded by later theorems.

\end{formaldefinition}

\begin{formaldefinition}[EasyCrypt line 96]
\noindent\textbf{EasyCrypt name.}
\begin{lstlisting}[language=EasyCrypt,basicstyle=\ttfamily\scriptsize]
main
\end{lstlisting}
\noindent\textbf{Checked EasyCrypt statement.}
\begin{lstlisting}[language=EasyCrypt,basicstyle=\ttfamily\scriptsize]
proc main(md : mode) : bool = {
\end{lstlisting}
\noindent\textbf{Formal mathematical reading.}
Mathematical definition. This procedure is a command in the probabilistic program.  Its return value, state updates, and oracle calls are the operational content of the game.

\noindent\textbf{Role in the proof.}
Use in this chapter. It supplies an executable component of the formal probabilistic model.

\end{formaldefinition}

\begin{formaldefinition}[EasyCrypt line 106]
\noindent\textbf{EasyCrypt name.}
\begin{lstlisting}[language=EasyCrypt,basicstyle=\ttfamily\scriptsize]
MLWE_Random_Game
\end{lstlisting}
\noindent\textbf{Checked EasyCrypt statement.}
\begin{lstlisting}[language=EasyCrypt,basicstyle=\ttfamily\scriptsize]
module MLWE_Random_Game(B : MLWE_Distinguisher) = {
\end{lstlisting}
\noindent\textbf{Formal mathematical reading.}
Mathematical definition. This is a probabilistic program/game or a module adapter.  Its procedures induce the probability space used by the game-based proof.

\noindent\textbf{Role in the proof.}
Use in this chapter. It defines the game or game interface whose success event is bounded by later theorems.

\end{formaldefinition}

\begin{formaldefinition}[EasyCrypt line 107]
\noindent\textbf{EasyCrypt name.}
\begin{lstlisting}[language=EasyCrypt,basicstyle=\ttfamily\scriptsize]
main
\end{lstlisting}
\noindent\textbf{Checked EasyCrypt statement.}
\begin{lstlisting}[language=EasyCrypt,basicstyle=\ttfamily\scriptsize]
proc main(md : mode) : bool = {
\end{lstlisting}
\noindent\textbf{Formal mathematical reading.}
Mathematical definition. This procedure is a command in the probabilistic program.  Its return value, state updates, and oracle calls are the operational content of the game.

\noindent\textbf{Role in the proof.}
Use in this chapter. It supplies an executable component of the formal probabilistic model.

\end{formaldefinition}

\begin{formaldefinition}[EasyCrypt line 117]
\noindent\textbf{EasyCrypt name.}
\begin{lstlisting}[language=EasyCrypt,basicstyle=\ttfamily\scriptsize]
ModuleSIS_Solver
\end{lstlisting}
\noindent\textbf{Checked EasyCrypt statement.}
\begin{lstlisting}[language=EasyCrypt,basicstyle=\ttfamily\scriptsize]
module type ModuleSIS_Solver = {
\end{lstlisting}
\noindent\textbf{Formal mathematical reading.}
Mathematical definition. This is an interface for an oracle, adversary, scheme, sampler, solver, or reduction.  A later theorem quantified over this interface is universally quantified over all modules implementing these procedures.

\noindent\textbf{Role in the proof.}
Use in this chapter. It supplies an executable component of the formal probabilistic model.

\end{formaldefinition}

\begin{formaldefinition}[EasyCrypt line 118]
\noindent\textbf{EasyCrypt name.}
\begin{lstlisting}[language=EasyCrypt,basicstyle=\ttfamily\scriptsize]
solve
\end{lstlisting}
\noindent\textbf{Checked EasyCrypt statement.}
\begin{lstlisting}[language=EasyCrypt,basicstyle=\ttfamily\scriptsize]
proc solve(x : module_sis_instance) : module_sis_solution option
}.
\end{lstlisting}
\noindent\textbf{Formal mathematical reading.}
Mathematical definition. This procedure is a command in the probabilistic program.  Its return value, state updates, and oracle calls are the operational content of the game.

\noindent\textbf{Role in the proof.}
Use in this chapter. It supplies an executable component of the formal probabilistic model.

\end{formaldefinition}

\begin{formaldefinition}[EasyCrypt line 121]
\noindent\textbf{EasyCrypt name.}
\begin{lstlisting}[language=EasyCrypt,basicstyle=\ttfamily\scriptsize]
ModuleSIS_Relation_Game
\end{lstlisting}
\noindent\textbf{Checked EasyCrypt statement.}
\begin{lstlisting}[language=EasyCrypt,basicstyle=\ttfamily\scriptsize]
module ModuleSIS_Relation_Game(B : ModuleSIS_Solver) = {
\end{lstlisting}
\noindent\textbf{Formal mathematical reading.}
Mathematical definition. This is a probabilistic program/game or a module adapter.  Its procedures induce the probability space used by the game-based proof.

\noindent\textbf{Role in the proof.}
Use in this chapter. It defines the game or game interface whose success event is bounded by later theorems.

\end{formaldefinition}

\begin{formaldefinition}[EasyCrypt line 122]
\noindent\textbf{EasyCrypt name.}
\begin{lstlisting}[language=EasyCrypt,basicstyle=\ttfamily\scriptsize]
main
\end{lstlisting}
\noindent\textbf{Checked EasyCrypt statement.}
\begin{lstlisting}[language=EasyCrypt,basicstyle=\ttfamily\scriptsize]
proc main(md : mode) : bool = {
\end{lstlisting}
\noindent\textbf{Formal mathematical reading.}
Mathematical definition. This procedure is a command in the probabilistic program.  Its return value, state updates, and oracle calls are the operational content of the game.

\noindent\textbf{Role in the proof.}
Use in this chapter. It supplies an executable component of the formal probabilistic model.

\end{formaldefinition}

\begin{formaldefinition}[EasyCrypt line 132]
\noindent\textbf{EasyCrypt name.}
\begin{lstlisting}[language=EasyCrypt,basicstyle=\ttfamily\scriptsize]
BimodalSelfTargetMSIS_Solver
\end{lstlisting}
\noindent\textbf{Checked EasyCrypt statement.}
\begin{lstlisting}[language=EasyCrypt,basicstyle=\ttfamily\scriptsize]
module type BimodalSelfTargetMSIS_Solver = {
\end{lstlisting}
\noindent\textbf{Formal mathematical reading.}
Mathematical definition. This is an interface for an oracle, adversary, scheme, sampler, solver, or reduction.  A later theorem quantified over this interface is universally quantified over all modules implementing these procedures.

\noindent\textbf{Role in the proof.}
Use in this chapter. It supplies an executable component of the formal probabilistic model.

\end{formaldefinition}

\begin{formaldefinition}[EasyCrypt line 133]
\noindent\textbf{EasyCrypt name.}
\begin{lstlisting}[language=EasyCrypt,basicstyle=\ttfamily\scriptsize]
solve
\end{lstlisting}
\noindent\textbf{Checked EasyCrypt statement.}
\begin{lstlisting}[language=EasyCrypt,basicstyle=\ttfamily\scriptsize]
proc solve(x : bimodal_selftarget_msis_instance) :
    bimodal_selftarget_msis_solution option
}.
\end{lstlisting}
\noindent\textbf{Formal mathematical reading.}
Mathematical definition. This procedure is a command in the probabilistic program.  Its return value, state updates, and oracle calls are the operational content of the game.

\noindent\textbf{Role in the proof.}
Use in this chapter. It supplies an executable component of the formal probabilistic model.

\end{formaldefinition}

\begin{formaldefinition}[EasyCrypt line 137]
\noindent\textbf{EasyCrypt name.}
\begin{lstlisting}[language=EasyCrypt,basicstyle=\ttfamily\scriptsize]
BimodalSelfTargetMSIS_Relation_Game
\end{lstlisting}
\noindent\textbf{Checked EasyCrypt statement.}
\begin{lstlisting}[language=EasyCrypt,basicstyle=\ttfamily\scriptsize]
module BimodalSelfTargetMSIS_Relation_Game(
  B : BimodalSelfTargetMSIS_Solver
) = {
\end{lstlisting}
\noindent\textbf{Formal mathematical reading.}
Mathematical definition. This is a probabilistic program/game or a module adapter.  Its procedures induce the probability space used by the game-based proof.

\noindent\textbf{Role in the proof.}
Use in this chapter. It defines the game or game interface whose success event is bounded by later theorems.

\end{formaldefinition}

\begin{formaldefinition}[EasyCrypt line 140]
\noindent\textbf{EasyCrypt name.}
\begin{lstlisting}[language=EasyCrypt,basicstyle=\ttfamily\scriptsize]
main
\end{lstlisting}
\noindent\textbf{Checked EasyCrypt statement.}
\begin{lstlisting}[language=EasyCrypt,basicstyle=\ttfamily\scriptsize]
proc main(md : mode) : bool = {
\end{lstlisting}
\noindent\textbf{Formal mathematical reading.}
Mathematical definition. This procedure is a command in the probabilistic program.  Its return value, state updates, and oracle calls are the operational content of the game.

\noindent\textbf{Role in the proof.}
Use in this chapter. It supplies an executable component of the formal probabilistic model.

\end{formaldefinition}

\begin{formaldefinition}[EasyCrypt line 150]
\noindent\textbf{EasyCrypt name.}
\begin{lstlisting}[language=EasyCrypt,basicstyle=\ttfamily\scriptsize]
lift_bimodal_solution_option
\end{lstlisting}
\noindent\textbf{Checked EasyCrypt statement.}
\begin{lstlisting}[language=EasyCrypt,basicstyle=\ttfamily\scriptsize]
op lift_bimodal_solution_option
   (md : mode) (x : bimodal_selftarget_msis_instance)
   (sol : bimodal_selftarget_msis_solution option) :
   module_sis_solution option =
  if sol = None then None
  else Some (bimodal_to_module_sis_solution md x (oget sol)).
\end{lstlisting}
\noindent\textbf{Formal mathematical reading.}
Mathematical definition. This is a total EasyCrypt operation.  The exact displayed statement gives the domain, codomain, and defining equation when present.

\noindent\textbf{Role in the proof.}
Use in this chapter. It is part of the exact formal vocabulary used by subsequent lemmas and games.

\end{formaldefinition}

\begin{formaldefinition}[EasyCrypt line 157]
\noindent\textbf{EasyCrypt name.}
\begin{lstlisting}[language=EasyCrypt,basicstyle=\ttfamily\scriptsize]
dmodule_sis_from_bimodal
\end{lstlisting}
\noindent\textbf{Checked EasyCrypt statement.}
\begin{lstlisting}[language=EasyCrypt,basicstyle=\ttfamily\scriptsize]
op dmodule_sis_from_bimodal (md : mode) : module_sis_instance distr =
  dmap (dbimodal_selftarget_msis_instance md)
    (bimodal_to_module_sis_instance md).
\end{lstlisting}
\noindent\textbf{Formal mathematical reading.}
Mathematical definition. This is a total EasyCrypt operation.  The exact displayed statement gives the domain, codomain, and defining equation when present.

\noindent\textbf{Role in the proof.}
Use in this chapter. It contributes a sampler or sampler-derived object to the distributional model.

\end{formaldefinition}

\begin{formaldefinition}[EasyCrypt line 174]
\noindent\textbf{EasyCrypt name.}
\begin{lstlisting}[language=EasyCrypt,basicstyle=\ttfamily\scriptsize]
BimodalToModuleSIS_Lift_Game
\end{lstlisting}
\noindent\textbf{Checked EasyCrypt statement.}
\begin{lstlisting}[language=EasyCrypt,basicstyle=\ttfamily\scriptsize]
module BimodalToModuleSIS_Lift_Game(
  B : BimodalSelfTargetMSIS_Solver
) = {
\end{lstlisting}
\noindent\textbf{Formal mathematical reading.}
Mathematical definition. This is a probabilistic program/game or a module adapter.  Its procedures induce the probability space used by the game-based proof.

\noindent\textbf{Role in the proof.}
Use in this chapter. It defines the game or game interface whose success event is bounded by later theorems.

\end{formaldefinition}

\begin{formaldefinition}[EasyCrypt line 177]
\noindent\textbf{EasyCrypt name.}
\begin{lstlisting}[language=EasyCrypt,basicstyle=\ttfamily\scriptsize]
main
\end{lstlisting}
\noindent\textbf{Checked EasyCrypt statement.}
\begin{lstlisting}[language=EasyCrypt,basicstyle=\ttfamily\scriptsize]
proc main(md : mode) : bool = {
\end{lstlisting}
\noindent\textbf{Formal mathematical reading.}
Mathematical definition. This procedure is a command in the probabilistic program.  Its return value, state updates, and oracle calls are the operational content of the game.

\noindent\textbf{Role in the proof.}
Use in this chapter. It supplies an executable component of the formal probabilistic model.

\end{formaldefinition}

\begin{formaldefinition}[EasyCrypt line 227]
\noindent\textbf{EasyCrypt name.}
\begin{lstlisting}[language=EasyCrypt,basicstyle=\ttfamily\scriptsize]
ModuleSIS_ModeSolver
\end{lstlisting}
\noindent\textbf{Checked EasyCrypt statement.}
\begin{lstlisting}[language=EasyCrypt,basicstyle=\ttfamily\scriptsize]
module type ModuleSIS_ModeSolver = {
\end{lstlisting}
\noindent\textbf{Formal mathematical reading.}
Mathematical definition. This is an interface for an oracle, adversary, scheme, sampler, solver, or reduction.  A later theorem quantified over this interface is universally quantified over all modules implementing these procedures.

\noindent\textbf{Role in the proof.}
Use in this chapter. It supplies an executable component of the formal probabilistic model.

\end{formaldefinition}

\begin{formaldefinition}[EasyCrypt line 228]
\noindent\textbf{EasyCrypt name.}
\begin{lstlisting}[language=EasyCrypt,basicstyle=\ttfamily\scriptsize]
solve
\end{lstlisting}
\noindent\textbf{Checked EasyCrypt statement.}
\begin{lstlisting}[language=EasyCrypt,basicstyle=\ttfamily\scriptsize]
proc solve(md : mode, x : module_sis_instance) :
    module_sis_solution option
}.
\end{lstlisting}
\noindent\textbf{Formal mathematical reading.}
Mathematical definition. This procedure is a command in the probabilistic program.  Its return value, state updates, and oracle calls are the operational content of the game.

\noindent\textbf{Role in the proof.}
Use in this chapter. It supplies an executable component of the formal probabilistic model.

\end{formaldefinition}

\begin{formaldefinition}[EasyCrypt line 232]
\noindent\textbf{EasyCrypt name.}
\begin{lstlisting}[language=EasyCrypt,basicstyle=\ttfamily\scriptsize]
BimodalSelfTargetMSIS_ModeSolver
\end{lstlisting}
\noindent\textbf{Checked EasyCrypt statement.}
\begin{lstlisting}[language=EasyCrypt,basicstyle=\ttfamily\scriptsize]
module type BimodalSelfTargetMSIS_ModeSolver = {
\end{lstlisting}
\noindent\textbf{Formal mathematical reading.}
Mathematical definition. This is an interface for an oracle, adversary, scheme, sampler, solver, or reduction.  A later theorem quantified over this interface is universally quantified over all modules implementing these procedures.

\noindent\textbf{Role in the proof.}
Use in this chapter. It supplies an executable component of the formal probabilistic model.

\end{formaldefinition}

\begin{formaldefinition}[EasyCrypt line 233]
\noindent\textbf{EasyCrypt name.}
\begin{lstlisting}[language=EasyCrypt,basicstyle=\ttfamily\scriptsize]
solve
\end{lstlisting}
\noindent\textbf{Checked EasyCrypt statement.}
\begin{lstlisting}[language=EasyCrypt,basicstyle=\ttfamily\scriptsize]
proc solve(md : mode, x : bimodal_selftarget_msis_instance) :
    bimodal_selftarget_msis_solution option
}.
\end{lstlisting}
\noindent\textbf{Formal mathematical reading.}
Mathematical definition. This procedure is a command in the probabilistic program.  Its return value, state updates, and oracle calls are the operational content of the game.

\noindent\textbf{Role in the proof.}
Use in this chapter. It supplies an executable component of the formal probabilistic model.

\end{formaldefinition}

\begin{formaldefinition}[EasyCrypt line 237]
\noindent\textbf{EasyCrypt name.}
\begin{lstlisting}[language=EasyCrypt,basicstyle=\ttfamily\scriptsize]
ModuleSIS_ModeRelation_Game
\end{lstlisting}
\noindent\textbf{Checked EasyCrypt statement.}
\begin{lstlisting}[language=EasyCrypt,basicstyle=\ttfamily\scriptsize]
module ModuleSIS_ModeRelation_Game(
  B : ModuleSIS_ModeSolver
) = {
\end{lstlisting}
\noindent\textbf{Formal mathematical reading.}
Mathematical definition. This is a probabilistic program/game or a module adapter.  Its procedures induce the probability space used by the game-based proof.

\noindent\textbf{Role in the proof.}
Use in this chapter. It defines the game or game interface whose success event is bounded by later theorems.

\end{formaldefinition}

\begin{formaldefinition}[EasyCrypt line 240]
\noindent\textbf{EasyCrypt name.}
\begin{lstlisting}[language=EasyCrypt,basicstyle=\ttfamily\scriptsize]
main
\end{lstlisting}
\noindent\textbf{Checked EasyCrypt statement.}
\begin{lstlisting}[language=EasyCrypt,basicstyle=\ttfamily\scriptsize]
proc main(md : mode) : bool = {
\end{lstlisting}
\noindent\textbf{Formal mathematical reading.}
Mathematical definition. This procedure is a command in the probabilistic program.  Its return value, state updates, and oracle calls are the operational content of the game.

\noindent\textbf{Role in the proof.}
Use in this chapter. It supplies an executable component of the formal probabilistic model.

\end{formaldefinition}

\begin{formaldefinition}[EasyCrypt line 250]
\noindent\textbf{EasyCrypt name.}
\begin{lstlisting}[language=EasyCrypt,basicstyle=\ttfamily\scriptsize]
BimodalSelfTargetMSIS_ModeRelation_Game
\end{lstlisting}
\noindent\textbf{Checked EasyCrypt statement.}
\begin{lstlisting}[language=EasyCrypt,basicstyle=\ttfamily\scriptsize]
module BimodalSelfTargetMSIS_ModeRelation_Game(
  B : BimodalSelfTargetMSIS_ModeSolver
) = {
\end{lstlisting}
\noindent\textbf{Formal mathematical reading.}
Mathematical definition. This is a probabilistic program/game or a module adapter.  Its procedures induce the probability space used by the game-based proof.

\noindent\textbf{Role in the proof.}
Use in this chapter. It defines the game or game interface whose success event is bounded by later theorems.

\end{formaldefinition}

\begin{formaldefinition}[EasyCrypt line 253]
\noindent\textbf{EasyCrypt name.}
\begin{lstlisting}[language=EasyCrypt,basicstyle=\ttfamily\scriptsize]
main
\end{lstlisting}
\noindent\textbf{Checked EasyCrypt statement.}
\begin{lstlisting}[language=EasyCrypt,basicstyle=\ttfamily\scriptsize]
proc main(md : mode) : bool = {
\end{lstlisting}
\noindent\textbf{Formal mathematical reading.}
Mathematical definition. This procedure is a command in the probabilistic program.  Its return value, state updates, and oracle calls are the operational content of the game.

\noindent\textbf{Role in the proof.}
Use in this chapter. It supplies an executable component of the formal probabilistic model.

\end{formaldefinition}

\begin{formaldefinition}[EasyCrypt line 263]
\noindent\textbf{EasyCrypt name.}
\begin{lstlisting}[language=EasyCrypt,basicstyle=\ttfamily\scriptsize]
ModuleSIS_LiftedDistribution_Relation_Game
\end{lstlisting}
\noindent\textbf{Checked EasyCrypt statement.}
\begin{lstlisting}[language=EasyCrypt,basicstyle=\ttfamily\scriptsize]
module ModuleSIS_LiftedDistribution_Relation_Game(
  B : ModuleSIS_ModeSolver
) = {
\end{lstlisting}
\noindent\textbf{Formal mathematical reading.}
Mathematical definition. This is a probabilistic program/game or a module adapter.  Its procedures induce the probability space used by the game-based proof.

\noindent\textbf{Role in the proof.}
Use in this chapter. It defines the game or game interface whose success event is bounded by later theorems.

\end{formaldefinition}

\begin{formaldefinition}[EasyCrypt line 266]
\noindent\textbf{EasyCrypt name.}
\begin{lstlisting}[language=EasyCrypt,basicstyle=\ttfamily\scriptsize]
main
\end{lstlisting}
\noindent\textbf{Checked EasyCrypt statement.}
\begin{lstlisting}[language=EasyCrypt,basicstyle=\ttfamily\scriptsize]
proc main(md : mode) : bool = {
\end{lstlisting}
\noindent\textbf{Formal mathematical reading.}
Mathematical definition. This procedure is a command in the probabilistic program.  Its return value, state updates, and oracle calls are the operational content of the game.

\noindent\textbf{Role in the proof.}
Use in this chapter. It supplies an executable component of the formal probabilistic model.

\end{formaldefinition}

\begin{formaldefinition}[EasyCrypt line 276]
\noindent\textbf{EasyCrypt name.}
\begin{lstlisting}[language=EasyCrypt,basicstyle=\ttfamily\scriptsize]
BimodalModeSolver_As_ModuleSIS
\end{lstlisting}
\noindent\textbf{Checked EasyCrypt statement.}
\begin{lstlisting}[language=EasyCrypt,basicstyle=\ttfamily\scriptsize]
module BimodalModeSolver_As_ModuleSIS(
  B : BimodalSelfTargetMSIS_ModeSolver
) = {
\end{lstlisting}
\noindent\textbf{Formal mathematical reading.}
Mathematical definition. This is a probabilistic program/game or a module adapter.  Its procedures induce the probability space used by the game-based proof.

\noindent\textbf{Role in the proof.}
Use in this chapter. It supplies an executable component of the formal probabilistic model.

\end{formaldefinition}

\begin{formaldefinition}[EasyCrypt line 279]
\noindent\textbf{EasyCrypt name.}
\begin{lstlisting}[language=EasyCrypt,basicstyle=\ttfamily\scriptsize]
solve
\end{lstlisting}
\noindent\textbf{Checked EasyCrypt statement.}
\begin{lstlisting}[language=EasyCrypt,basicstyle=\ttfamily\scriptsize]
proc solve(md : mode, x : module_sis_instance) :
    module_sis_solution option = {
\end{lstlisting}
\noindent\textbf{Formal mathematical reading.}
Mathematical definition. This procedure is a command in the probabilistic program.  Its return value, state updates, and oracle calls are the operational content of the game.

\noindent\textbf{Role in the proof.}
Use in this chapter. It supplies an executable component of the formal probabilistic model.

\end{formaldefinition}

\begin{formaldefinition}[EasyCrypt line 310]
\noindent\textbf{EasyCrypt name.}
\begin{lstlisting}[language=EasyCrypt,basicstyle=\ttfamily\scriptsize]
MLWE_Reduction
\end{lstlisting}
\noindent\textbf{Checked EasyCrypt statement.}
\begin{lstlisting}[language=EasyCrypt,basicstyle=\ttfamily\scriptsize]
module type MLWE_Reduction = {
\end{lstlisting}
\noindent\textbf{Formal mathematical reading.}
Mathematical definition. This is an interface for an oracle, adversary, scheme, sampler, solver, or reduction.  A later theorem quantified over this interface is universally quantified over all modules implementing these procedures.

\noindent\textbf{Role in the proof.}
Use in this chapter. It supplies an executable component of the formal probabilistic model.

\end{formaldefinition}

\begin{formaldefinition}[EasyCrypt line 311]
\noindent\textbf{EasyCrypt name.}
\begin{lstlisting}[language=EasyCrypt,basicstyle=\ttfamily\scriptsize]
main
\end{lstlisting}
\noindent\textbf{Checked EasyCrypt statement.}
\begin{lstlisting}[language=EasyCrypt,basicstyle=\ttfamily\scriptsize]
proc main() : bool
}.
\end{lstlisting}
\noindent\textbf{Formal mathematical reading.}
Mathematical definition. This procedure is a command in the probabilistic program.  Its return value, state updates, and oracle calls are the operational content of the game.

\noindent\textbf{Role in the proof.}
Use in this chapter. It supplies an executable component of the formal probabilistic model.

\end{formaldefinition}

\begin{formaldefinition}[EasyCrypt line 314]
\noindent\textbf{EasyCrypt name.}
\begin{lstlisting}[language=EasyCrypt,basicstyle=\ttfamily\scriptsize]
BimodalSelfTargetMSIS_Reduction
\end{lstlisting}
\noindent\textbf{Checked EasyCrypt statement.}
\begin{lstlisting}[language=EasyCrypt,basicstyle=\ttfamily\scriptsize]
module type BimodalSelfTargetMSIS_Reduction = {
\end{lstlisting}
\noindent\textbf{Formal mathematical reading.}
Mathematical definition. This is an interface for an oracle, adversary, scheme, sampler, solver, or reduction.  A later theorem quantified over this interface is universally quantified over all modules implementing these procedures.

\noindent\textbf{Role in the proof.}
Use in this chapter. It supplies an executable component of the formal probabilistic model.

\end{formaldefinition}

\begin{formaldefinition}[EasyCrypt line 315]
\noindent\textbf{EasyCrypt name.}
\begin{lstlisting}[language=EasyCrypt,basicstyle=\ttfamily\scriptsize]
main
\end{lstlisting}
\noindent\textbf{Checked EasyCrypt statement.}
\begin{lstlisting}[language=EasyCrypt,basicstyle=\ttfamily\scriptsize]
proc main() : bool
}.
\end{lstlisting}
\noindent\textbf{Formal mathematical reading.}
Mathematical definition. This procedure is a command in the probabilistic program.  Its return value, state updates, and oracle calls are the operational content of the game.

\noindent\textbf{Role in the proof.}
Use in this chapter. It supplies an executable component of the formal probabilistic model.

\end{formaldefinition}

\begin{formaldefinition}[EasyCrypt line 318]
\noindent\textbf{EasyCrypt name.}
\begin{lstlisting}[language=EasyCrypt,basicstyle=\ttfamily\scriptsize]
ModuleSIS_Reduction
\end{lstlisting}
\noindent\textbf{Checked EasyCrypt statement.}
\begin{lstlisting}[language=EasyCrypt,basicstyle=\ttfamily\scriptsize]
module type ModuleSIS_Reduction = {
\end{lstlisting}
\noindent\textbf{Formal mathematical reading.}
Mathematical definition. This is an interface for an oracle, adversary, scheme, sampler, solver, or reduction.  A later theorem quantified over this interface is universally quantified over all modules implementing these procedures.

\noindent\textbf{Role in the proof.}
Use in this chapter. It supplies an executable component of the formal probabilistic model.

\end{formaldefinition}

\begin{formaldefinition}[EasyCrypt line 319]
\noindent\textbf{EasyCrypt name.}
\begin{lstlisting}[language=EasyCrypt,basicstyle=\ttfamily\scriptsize]
main
\end{lstlisting}
\noindent\textbf{Checked EasyCrypt statement.}
\begin{lstlisting}[language=EasyCrypt,basicstyle=\ttfamily\scriptsize]
proc main() : bool
}.
\end{lstlisting}
\noindent\textbf{Formal mathematical reading.}
Mathematical definition. This procedure is a command in the probabilistic program.  Its return value, state updates, and oracle calls are the operational content of the game.

\noindent\textbf{Role in the proof.}
Use in this chapter. It supplies an executable component of the formal probabilistic model.

\end{formaldefinition}

\begin{formaldefinition}[EasyCrypt line 322]
\noindent\textbf{EasyCrypt name.}
\begin{lstlisting}[language=EasyCrypt,basicstyle=\ttfamily\scriptsize]
MLWE_Game
\end{lstlisting}
\noindent\textbf{Checked EasyCrypt statement.}
\begin{lstlisting}[language=EasyCrypt,basicstyle=\ttfamily\scriptsize]
module MLWE_Game(B : MLWE_Reduction) = {
\end{lstlisting}
\noindent\textbf{Formal mathematical reading.}
Mathematical definition. This is a probabilistic program/game or a module adapter.  Its procedures induce the probability space used by the game-based proof.

\noindent\textbf{Role in the proof.}
Use in this chapter. It defines the game or game interface whose success event is bounded by later theorems.

\end{formaldefinition}

\begin{formaldefinition}[EasyCrypt line 323]
\noindent\textbf{EasyCrypt name.}
\begin{lstlisting}[language=EasyCrypt,basicstyle=\ttfamily\scriptsize]
main
\end{lstlisting}
\noindent\textbf{Checked EasyCrypt statement.}
\begin{lstlisting}[language=EasyCrypt,basicstyle=\ttfamily\scriptsize]
proc main() : bool = {
\end{lstlisting}
\noindent\textbf{Formal mathematical reading.}
Mathematical definition. This procedure is a command in the probabilistic program.  Its return value, state updates, and oracle calls are the operational content of the game.

\noindent\textbf{Role in the proof.}
Use in this chapter. It supplies an executable component of the formal probabilistic model.

\end{formaldefinition}

\begin{formaldefinition}[EasyCrypt line 331]
\noindent\textbf{EasyCrypt name.}
\begin{lstlisting}[language=EasyCrypt,basicstyle=\ttfamily\scriptsize]
BimodalSelfTargetMSIS_Game
\end{lstlisting}
\noindent\textbf{Checked EasyCrypt statement.}
\begin{lstlisting}[language=EasyCrypt,basicstyle=\ttfamily\scriptsize]
module BimodalSelfTargetMSIS_Game(B : BimodalSelfTargetMSIS_Reduction) = {
\end{lstlisting}
\noindent\textbf{Formal mathematical reading.}
Mathematical definition. This is a probabilistic program/game or a module adapter.  Its procedures induce the probability space used by the game-based proof.

\noindent\textbf{Role in the proof.}
Use in this chapter. It defines the game or game interface whose success event is bounded by later theorems.

\end{formaldefinition}

\begin{formaldefinition}[EasyCrypt line 332]
\noindent\textbf{EasyCrypt name.}
\begin{lstlisting}[language=EasyCrypt,basicstyle=\ttfamily\scriptsize]
main
\end{lstlisting}
\noindent\textbf{Checked EasyCrypt statement.}
\begin{lstlisting}[language=EasyCrypt,basicstyle=\ttfamily\scriptsize]
proc main() : bool = {
\end{lstlisting}
\noindent\textbf{Formal mathematical reading.}
Mathematical definition. This procedure is a command in the probabilistic program.  Its return value, state updates, and oracle calls are the operational content of the game.

\noindent\textbf{Role in the proof.}
Use in this chapter. It supplies an executable component of the formal probabilistic model.

\end{formaldefinition}

\begin{formaldefinition}[EasyCrypt line 340]
\noindent\textbf{EasyCrypt name.}
\begin{lstlisting}[language=EasyCrypt,basicstyle=\ttfamily\scriptsize]
ModuleSIS_Game
\end{lstlisting}
\noindent\textbf{Checked EasyCrypt statement.}
\begin{lstlisting}[language=EasyCrypt,basicstyle=\ttfamily\scriptsize]
module ModuleSIS_Game(B : ModuleSIS_Reduction) = {
\end{lstlisting}
\noindent\textbf{Formal mathematical reading.}
Mathematical definition. This is a probabilistic program/game or a module adapter.  Its procedures induce the probability space used by the game-based proof.

\noindent\textbf{Role in the proof.}
Use in this chapter. It defines the game or game interface whose success event is bounded by later theorems.

\end{formaldefinition}

\begin{formaldefinition}[EasyCrypt line 341]
\noindent\textbf{EasyCrypt name.}
\begin{lstlisting}[language=EasyCrypt,basicstyle=\ttfamily\scriptsize]
main
\end{lstlisting}
\noindent\textbf{Checked EasyCrypt statement.}
\begin{lstlisting}[language=EasyCrypt,basicstyle=\ttfamily\scriptsize]
proc main() : bool = {
\end{lstlisting}
\noindent\textbf{Formal mathematical reading.}
Mathematical definition. This procedure is a command in the probabilistic program.  Its return value, state updates, and oracle calls are the operational content of the game.

\noindent\textbf{Role in the proof.}
Use in this chapter. It supplies an executable component of the formal probabilistic model.

\end{formaldefinition}

\begin{formaldefinition}[EasyCrypt line 349]
\noindent\textbf{EasyCrypt name.}
\begin{lstlisting}[language=EasyCrypt,basicstyle=\ttfamily\scriptsize]
BimodalMSIS_As_ModuleSIS
\end{lstlisting}
\noindent\textbf{Checked EasyCrypt statement.}
\begin{lstlisting}[language=EasyCrypt,basicstyle=\ttfamily\scriptsize]
module BimodalMSIS_As_ModuleSIS(B : BimodalSelfTargetMSIS_Reduction) = {
\end{lstlisting}
\noindent\textbf{Formal mathematical reading.}
Mathematical definition. This is a probabilistic program/game or a module adapter.  Its procedures induce the probability space used by the game-based proof.

\noindent\textbf{Role in the proof.}
Use in this chapter. It supplies an executable component of the formal probabilistic model.

\end{formaldefinition}

\begin{formaldefinition}[EasyCrypt line 350]
\noindent\textbf{EasyCrypt name.}
\begin{lstlisting}[language=EasyCrypt,basicstyle=\ttfamily\scriptsize]
main
\end{lstlisting}
\noindent\textbf{Checked EasyCrypt statement.}
\begin{lstlisting}[language=EasyCrypt,basicstyle=\ttfamily\scriptsize]
proc main() : bool = {
\end{lstlisting}
\noindent\textbf{Formal mathematical reading.}
Mathematical definition. This procedure is a command in the probabilistic program.  Its return value, state updates, and oracle calls are the operational content of the game.

\noindent\textbf{Role in the proof.}
Use in this chapter. It supplies an executable component of the formal probabilistic model.

\end{formaldefinition}

\subsubsection*{Lemmas, Theorems, Relational Equivalences, and Their Proof Dependencies}
\begin{formaltheorem}[EasyCrypt line 43]
\noindent\textbf{EasyCrypt name.}
\begin{lstlisting}[language=EasyCrypt,basicstyle=\ttfamily\scriptsize]
module_sis_eval_zero_instance
\end{lstlisting}
\noindent\textbf{Checked EasyCrypt statement.}
\begin{lstlisting}[language=EasyCrypt,basicstyle=\ttfamily\scriptsize]
lemma module_sis_eval_zero_instance md sol :
  module_sis_eval md (module_sis_zero_instance md) sol = polyveck_zero md.
\end{lstlisting}
\noindent\textbf{Formal mathematical reading.}
Mathematical theorem. This is an equality or definitional expansion used for rewriting and normalization.

\noindent\textbf{Role in the proof.}
Use in this chapter. It contributes directly to an advantage bound, bad-event bound, or real-arithmetic composition step.

\noindent\textbf{Proof-script interpretation.}
Proof-script reading. The machine proof decomposes the theorem into rewrites, program-logic obligations, relational equivalences, and arithmetic side conditions.
\emph{Main tactics visible in the proof script:} \ec{rewrite}.
\emph{Named identifiers syntactically cited in the proof block:}
\begin{lstlisting}[language=EasyCrypt,basicstyle=\ttfamily\scriptsize]
matrix_vec_mul
module_sis_eval
module_sis_zero_instance
polyveck_add
polyveck_neg
polyveck_sub
\end{lstlisting}

\end{formaltheorem}

\begin{formaltheorem}[EasyCrypt line 50]
\noindent\textbf{EasyCrypt name.}
\begin{lstlisting}[language=EasyCrypt,basicstyle=\ttfamily\scriptsize]
module_sis_zero_valid
\end{lstlisting}
\noindent\textbf{Checked EasyCrypt statement.}
\begin{lstlisting}[language=EasyCrypt,basicstyle=\ttfamily\scriptsize]
lemma module_sis_zero_valid md :
  module_sis_valid md (module_sis_zero_instance md) (polyvecl_zero md).
\end{lstlisting}
\noindent\textbf{Formal mathematical reading.}
Mathematical theorem. This lemma is a universally quantified proposition over the variables shown in the EasyCrypt statement.

\noindent\textbf{Role in the proof.}
Use in this chapter. It contributes directly to an advantage bound, bad-event bound, or real-arithmetic composition step.

\noindent\textbf{Proof-script interpretation.}
Proof-script reading. The machine proof decomposes the theorem into rewrites, program-logic obligations, relational equivalences, and arithmetic side conditions.
\emph{Main tactics visible in the proof script:} \ec{rewrite}, \ec{split}, \ec{apply}.
\emph{Named identifiers syntactically cited in the proof block:}
\begin{lstlisting}[language=EasyCrypt,basicstyle=\ttfamily\scriptsize]
matrix_vec_mul
module_sis_solution_wf
module_sis_valid
module_sis_zero_instance
polyvecl_zero_wf
\end{lstlisting}

\end{formaltheorem}

\begin{formaltheorem}[EasyCrypt line 59]
\noindent\textbf{EasyCrypt name.}
\begin{lstlisting}[language=EasyCrypt,basicstyle=\ttfamily\scriptsize]
module_sis_eval_wf
\end{lstlisting}
\noindent\textbf{Checked EasyCrypt statement.}
\begin{lstlisting}[language=EasyCrypt,basicstyle=\ttfamily\scriptsize]
lemma module_sis_eval_wf md x sol :
  polyveck_wf md (module_sis_eval md x sol).
\end{lstlisting}
\noindent\textbf{Formal mathematical reading.}
Mathematical theorem. This lemma is a universally quantified proposition over the variables shown in the EasyCrypt statement.

\noindent\textbf{Role in the proof.}
Use in this chapter. It contributes directly to an advantage bound, bad-event bound, or real-arithmetic composition step.

\noindent\textbf{Proof-script interpretation.}
No proof block was collected for this item.

\end{formaltheorem}

\begin{formaltheorem}[EasyCrypt line 63]
\noindent\textbf{EasyCrypt name.}
\begin{lstlisting}[language=EasyCrypt,basicstyle=\ttfamily\scriptsize]
module_sis_valid_target_wf
\end{lstlisting}
\noindent\textbf{Checked EasyCrypt statement.}
\begin{lstlisting}[language=EasyCrypt,basicstyle=\ttfamily\scriptsize]
lemma module_sis_valid_target_wf md x sol :
  module_sis_valid md x sol =>
  polyveck_wf md x.`2.
\end{lstlisting}
\noindent\textbf{Formal mathematical reading.}
Mathematical theorem. This is an implication, normally turning one invariant, validity predicate, or proof obligation into another.

\noindent\textbf{Role in the proof.}
Use in this chapter. It contributes directly to an advantage bound, bad-event bound, or real-arithmetic composition step.

\noindent\textbf{Proof-script interpretation.}
Proof-script reading. The machine proof decomposes the theorem into rewrites, program-logic obligations, relational equivalences, and arithmetic side conditions.
\emph{Main tactics visible in the proof script:} \ec{rewrite}, \ec{apply}.
\emph{Named identifiers syntactically cited in the proof block:}
\begin{lstlisting}[language=EasyCrypt,basicstyle=\ttfamily\scriptsize]
matrix_vec_mul_wf
module_sis_valid
targetE
\end{lstlisting}

\end{formaltheorem}

\begin{formaltheorem}[EasyCrypt line 80]
\noindent\textbf{EasyCrypt name.}
\begin{lstlisting}[language=EasyCrypt,basicstyle=\ttfamily\scriptsize]
bimodal_to_module_sis_valid
\end{lstlisting}
\noindent\textbf{Checked EasyCrypt statement.}
\begin{lstlisting}[language=EasyCrypt,basicstyle=\ttfamily\scriptsize]
lemma bimodal_to_module_sis_valid md x sol :
  bimodal_selftarget_msis_valid md x sol =>
  module_sis_valid md
    (bimodal_to_module_sis_instance md x)
    (bimodal_to_module_sis_solution md x sol).
\end{lstlisting}
\noindent\textbf{Formal mathematical reading.}
Mathematical theorem. This is an implication, normally turning one invariant, validity predicate, or proof obligation into another.

\noindent\textbf{Role in the proof.}
Use in this chapter. It contributes directly to an advantage bound, bad-event bound, or real-arithmetic composition step.

\noindent\textbf{Proof-script interpretation.}
Proof-script reading. The machine proof decomposes the theorem into rewrites, program-logic obligations, relational equivalences, and arithmetic side conditions.
\emph{Main tactics visible in the proof script:} \ec{rewrite}.
\emph{Named identifiers syntactically cited in the proof block:}
\begin{lstlisting}[language=EasyCrypt,basicstyle=\ttfamily\scriptsize]
bimodal_selftarget_msis_valid
bimodal_to_module_sis_instance
bimodal_to_module_sis_solution
\end{lstlisting}

\end{formaltheorem}

\begin{formaltheorem}[EasyCrypt line 161]
\noindent\textbf{EasyCrypt name.}
\begin{lstlisting}[language=EasyCrypt,basicstyle=\ttfamily\scriptsize]
dmodule_sis_from_bimodalE
\end{lstlisting}
\noindent\textbf{Checked EasyCrypt statement.}
\begin{lstlisting}[language=EasyCrypt,basicstyle=\ttfamily\scriptsize]
lemma dmodule_sis_from_bimodalE md :
  dmodule_sis_from_bimodal md = dbimodal_selftarget_msis_instance md.
\end{lstlisting}
\noindent\textbf{Formal mathematical reading.}
Mathematical theorem. This is an equality or definitional expansion used for rewriting and normalization.

\noindent\textbf{Role in the proof.}
Use in this chapter. It contributes directly to an advantage bound, bad-event bound, or real-arithmetic composition step.

\noindent\textbf{Proof-script interpretation.}
Proof-script reading. The machine proof decomposes the theorem into rewrites, program-logic obligations, relational equivalences, and arithmetic side conditions.
\emph{Main tactics visible in the proof script:} \ec{rewrite}.
\emph{Named identifiers syntactically cited in the proof block:}
\begin{lstlisting}[language=EasyCrypt,basicstyle=\ttfamily\scriptsize]
bimodal_to_module_sis_instance
dmap_id
dmodule_sis_from_bimodal
\end{lstlisting}

\end{formaltheorem}

\begin{formaltheorem}[EasyCrypt line 167]
\noindent\textbf{EasyCrypt name.}
\begin{lstlisting}[language=EasyCrypt,basicstyle=\ttfamily\scriptsize]
dmodule_sis_from_bimodal_lossless
\end{lstlisting}
\noindent\textbf{Checked EasyCrypt statement.}
\begin{lstlisting}[language=EasyCrypt,basicstyle=\ttfamily\scriptsize]
lemma dmodule_sis_from_bimodal_lossless md :
  is_lossless (dmodule_sis_from_bimodal md).
\end{lstlisting}
\noindent\textbf{Formal mathematical reading.}
Mathematical theorem. This lemma is a universally quantified proposition over the variables shown in the EasyCrypt statement.

\noindent\textbf{Role in the proof.}
Use in this chapter. It contributes directly to an advantage bound, bad-event bound, or real-arithmetic composition step.

\noindent\textbf{Proof-script interpretation.}
Proof-script reading. The machine proof decomposes the theorem into rewrites, program-logic obligations, relational equivalences, and arithmetic side conditions.
\emph{Main tactics visible in the proof script:} \ec{rewrite}, \ec{apply}.
\emph{Named identifiers syntactically cited in the proof block:}
\begin{lstlisting}[language=EasyCrypt,basicstyle=\ttfamily\scriptsize]
dbimodal_selftarget_msis_instance_ll
dmodule_sis_from_bimodalE
\end{lstlisting}

\end{formaltheorem}

\begin{formaltheorem}[EasyCrypt line 191]
\noindent\textbf{EasyCrypt name.}
\begin{lstlisting}[language=EasyCrypt,basicstyle=\ttfamily\scriptsize]
bimodal_lift_successE
\end{lstlisting}
\noindent\textbf{Checked EasyCrypt statement.}
\begin{lstlisting}[language=EasyCrypt,basicstyle=\ttfamily\scriptsize]
lemma bimodal_lift_successE md x sol :
  (sol <> None /\ bimodal_selftarget_msis_valid md x (oget sol)) =
  (lift_bimodal_solution_option md x sol <> None /\
   module_sis_valid md
     (bimodal_to_module_sis_instance md x)
     (oget (lift_bimodal_solution_option md x sol))).
\end{lstlisting}
\noindent\textbf{Formal mathematical reading.}
Mathematical theorem. This is an equality or definitional expansion used for rewriting and normalization.

\noindent\textbf{Role in the proof.}
Use in this chapter. It is a reusable checked theorem cited by later proof scripts.

\noindent\textbf{Proof-script interpretation.}
Proof-script reading. The machine proof decomposes the theorem into rewrites, program-logic obligations, relational equivalences, and arithmetic side conditions.
\emph{Main tactics visible in the proof script:} \ec{case}, \ec{rewrite}.
\emph{Named identifiers syntactically cited in the proof block:}
\begin{lstlisting}[language=EasyCrypt,basicstyle=\ttfamily\scriptsize]
None
bimodal_selftarget_msis_valid
bimodal_to_module_sis_instance
bimodal_to_module_sis_solution
lift_bimodal_solution_option
sol
sol_none
\end{lstlisting}

\end{formaltheorem}

\begin{formaltheorem}[EasyCrypt line 209]
\noindent\textbf{EasyCrypt name.}
\begin{lstlisting}[language=EasyCrypt,basicstyle=\ttfamily\scriptsize]
bimodal_solver_lift_exact
\end{lstlisting}
\noindent\textbf{Checked EasyCrypt statement.}
\begin{lstlisting}[language=EasyCrypt,basicstyle=\ttfamily\scriptsize]
lemma bimodal_solver_lift_exact &m md :
  Pr[BimodalSelfTargetMSIS_Relation_Game(B).main(md) @ &m : res] =
  Pr[BimodalToModuleSIS_Lift_Game(B).main(md) @ &m : res].
\end{lstlisting}
\noindent\textbf{Formal mathematical reading.}
Mathematical theorem. This theorem has the form \(\Pr[G:E] = B\) is a game-probability statement.  It is one of the exact hops, bad-event bounds, or final advantage inequalities used in the reduction.

\noindent\textbf{Role in the proof.}
Use in this chapter. It proves an exact game replacement or simulator equivalence.

\noindent\textbf{Proof-script interpretation.}
Proof-script reading. The machine proof decomposes the theorem into rewrites, program-logic obligations, relational equivalences, and arithmetic side conditions.
\emph{Main tactics visible in the proof script:} \ec{byequiv}, \ec{proc}, \ec{wp}, \ec{call}, \ec{rnd}, \ec{rewrite}.
\emph{Named identifiers syntactically cited in the proof block:}
\begin{lstlisting}[language=EasyCrypt,basicstyle=\ttfamily\scriptsize]
bimodal_selftarget_msis_valid
bimodal_to_module_sis_instance
bimodal_to_module_sis_solution
glob
md
\end{lstlisting}

\end{formaltheorem}

\begin{formaltheorem}[EasyCrypt line 292]
\noindent\textbf{EasyCrypt name.}
\begin{lstlisting}[language=EasyCrypt,basicstyle=\ttfamily\scriptsize]
bimodal_mode_solver_lifted_distribution_exact
\end{lstlisting}
\noindent\textbf{Checked EasyCrypt statement.}
\begin{lstlisting}[language=EasyCrypt,basicstyle=\ttfamily\scriptsize]
lemma bimodal_mode_solver_lifted_distribution_exact &m md :
  Pr[BimodalSelfTargetMSIS_ModeRelation_Game(B).main(md) @ &m : res] =
  Pr[ModuleSIS_LiftedDistribution_Relation_Game(
       BimodalModeSolver_As_ModuleSIS(B)).main(md) @ &m : res].
\end{lstlisting}
\noindent\textbf{Formal mathematical reading.}
Mathematical theorem. This theorem has the form \(\Pr[G:E] = B\) is a game-probability statement.  It is one of the exact hops, bad-event bounds, or final advantage inequalities used in the reduction.

\noindent\textbf{Role in the proof.}
Use in this chapter. It proves an exact game replacement or simulator equivalence.

\noindent\textbf{Proof-script interpretation.}
Proof-script reading. The machine proof decomposes the theorem into rewrites, program-logic obligations, relational equivalences, and arithmetic side conditions.
\emph{Main tactics visible in the proof script:} \ec{byequiv}, \ec{proc}, \ec{inline}, \ec{wp}, \ec{call}, \ec{rnd}, \ec{rewrite}.
\emph{Named identifiers syntactically cited in the proof block:}
\begin{lstlisting}[language=EasyCrypt,basicstyle=\ttfamily\scriptsize]
BimodalModeSolver_As_ModuleSIS
bimodal_selftarget_msis_valid
glob
md
solve
\end{lstlisting}

\end{formaltheorem}

\begin{formaltheorem}[EasyCrypt line 362]
\noindent\textbf{EasyCrypt name.}
\begin{lstlisting}[language=EasyCrypt,basicstyle=\ttfamily\scriptsize]
bimodal_msis_as_module_sis_exact
\end{lstlisting}
\noindent\textbf{Checked EasyCrypt statement.}
\begin{lstlisting}[language=EasyCrypt,basicstyle=\ttfamily\scriptsize]
lemma bimodal_msis_as_module_sis_exact &m :
  Pr[BimodalSelfTargetMSIS_Game(B).main() @ &m : res] =
  Pr[ModuleSIS_Game(BimodalMSIS_As_ModuleSIS(B)).main() @ &m : res].
\end{lstlisting}
\noindent\textbf{Formal mathematical reading.}
Mathematical theorem. This theorem has the form \(\Pr[G:E] = B\) is a game-probability statement.  It is one of the exact hops, bad-event bounds, or final advantage inequalities used in the reduction.

\noindent\textbf{Role in the proof.}
Use in this chapter. It contributes directly to an advantage bound, bad-event bound, or real-arithmetic composition step.

\noindent\textbf{Proof-script interpretation.}
Proof-script reading. The machine proof decomposes the theorem into rewrites, program-logic obligations, relational equivalences, and arithmetic side conditions.
\emph{Main tactics visible in the proof script:} \ec{byequiv}, \ec{proc}, \ec{inline}, \ec{wp}, \ec{call}.
\emph{Named identifiers syntactically cited in the proof block:}
\begin{lstlisting}[language=EasyCrypt,basicstyle=\ttfamily\scriptsize]
BimodalMSIS_As_ModuleSIS
glob
main
\end{lstlisting}

\end{formaltheorem}

\medskip
\noindent\emph{End of generated formal inventory for \file{HAETAE_Assumptions.ec}.}
